\section{SPI on four wires}
\label{sec:spilines}

\begin{booktable}{@{}llL@{}}
\thead{Wire} & \thead{Direction} & \thead{Meaning} \\ \midrule
\code{SCK} & master $\rightarrow$ slave & Clock. Nothing happens without edges on it. \\
\code{MOSI} & master $\rightarrow$ slave & The master's data. \\
\code{MISO} & slave $\rightarrow$ master & The slave's data. \\
\code{SS} & master $\rightarrow$ slave & Active low. Frames the transaction. \\
\end{booktable}

The most important property: \textbf{SPI is full duplex.} Every clock pulse simultaneously shifts
one bit out on \code{MOSI} and one bit in on \code{MISO}. There is no way to "just read" --- to
get a byte in you have to send a byte out.

That is why a read sends four dummy bytes of \code{0x00}, and why the four \code{MISO} bytes
during a write are meaningless junk.

The parameters: SPI mode 0 (the clock idles low, data is sampled on the rising edge), MSB first,
\code{SCK} at most 1~MHz.

\section{The transaction}
\label{sec:transaction}

\textbf{Every transaction is exactly five bytes}: one command byte followed by four data bytes,
with \code{SS} low the whole way.

\begin{console}
Bit:      7    6    5    4    3    2    1    0
        +----+----+----+----+----+----+----+----+
        | W  | 0  | 0  | 0  |   register index  |
        +----+----+----+----+----+----+----+----+
\end{console}

\code{W = 1} means a write, \code{W = 0} a read. Bits 6--4 are reserved and are to be \code{0}.
Bits 3--0 are the register index, 0--12, from the register map in \kapref{ch:can}.

\subsection{Writing}

The master sends the register value in the four data bytes, \textbf{MSB first}. The slave commits
the assembled value \textbf{only once the fifth byte is complete}.

\subsection{Reading}

At the end of the command byte the slave \textbf{latches the register's current value once}, into
its response shifter. The four data bytes then clock that value out.

\textbf{The latching rule is part of the contract, not an implementation detail.} The four bytes
always belong to \emph{one} coherent sample of the register, taken in a single instant. Without
that rule a \code{STATUS} read could have been stitched together from two points in time --- the
TX-ready bit from before a transmission and the RX-valid bit from after. That would have been a
bug appearing about once in a thousand times and never reproducible.

\subsection{The \texorpdfstring{\code{SS}}{SS} rules}

Three rules, and all three have cost somebody an evening:

\begin{itemize}
  \item \textbf{One transaction per low \code{SS} period.} The master has to let \code{SS} go high
    between transactions.
  \item \textbf{Back-to-back does not work, and fails silently.} A sixth byte arriving while
    \code{SS} is still low is \emph{discarded} --- it is not interpreted as a new command byte.
    The slave being idle is therefore not enough.
  \item \textbf{An abort has no side effects.} If \code{SS} goes high before the fifth byte, the
    transaction is abandoned entirely: nothing is committed, no trigger fires. That is what makes
    the slave self-healing after a glitch.
\end{itemize}

That is why \code{ByteTransport} in \kapref{ch:byte} has \code{select()} and \code{deselect()} as
methods of their own: the framing is the driver's decision, not the transport's.

\subsection{\texorpdfstring{\code{MISO}}{MISO} outside a read}

The byte the master clocks out \textbf{during the command byte} is meaningless. The slave loads
its response shifter at the \emph{end} of the command byte, so what leaves during it is leftovers
from last time. \textbf{Discard it.}

\subsection{A complete example}

Sending a frame (\code{id 0x123}, \code{dlc 2}, data \code{AA BB}):

\begin{console}
Write TX_ID      : 81 00 00 01 23
Write TX_DLC     : 82 00 00 00 02
Write TX_DATA_LO : 83 AA BB 00 00
Write TX_DATA_HI : 84 00 00 00 00
Write TX_SEND    : 85 00 00 00 01
Read  STATUS     : 00 xx xx xx xx   -> MISO gives 00 00 00 00 during transmission,
                                      then 00 00 00 01 when the port is free again.
\end{console}

Read off the command byte in every line: \code{0x81} is the write bit plus index 1, that is,
\code{TX_ID}. The \code{0x00} on the last line is a read of index 0.

Note \code{83 AA BB 00 00}: \textbf{data byte 0 in the most significant position}. It is the same
packing on both sides of the wire.

\section{What it costs in time}
\label{sec:timing}

A transaction is 5 bytes = 40 bits. At 1~MHz the clocking takes 40~µs, plus \code{SS} handling.
Reckon on \textbf{40--50~µs per transaction}:

\begin{narrowtable}{@{}lrl@{}}
\thead{Operation} & \thead{Transactions} & \thead{Roughly} \\ \midrule
\code{hasError()} & 1 & 45~µs \\
\code{receive()}, no frame & 1 & 45~µs \\
\code{receive()}, one frame & 6 & 270~µs \\
\code{send()} & 6 & 270~µs \\
\end{narrowtable}

Set that against the CAN side: a maximal CAN frame is around 130 bits, that is, \textbf{130~µs} at
1~Mbit/s.

\begin{limitation}
\textbf{A \code{send()} takes longer than the frame it sends.} The SPI link, not the CAN bus, is
the bottleneck in this design.

\textbf{And the polling can miss frames.} The controller does not buffer: if two frames arrive
close together and your \code{receive()} takes 270~µs, the second has time to overwrite the first.
That is the hardware's documented limitation, not your bug --- but it is your code that notices
it.
\end{limitation}
